\section{The empirical picture}

A common microscopic construction ties the galactic acceleration scale,
the initial gravitating stock, and its conversion into vacuum energy.
The tests below examine those connections across different formation
histories and established gravitational and cosmological observations.

Distinctive tests compare common coefficients and inherited stock with
specified alternatives. LMC masses, Crater II stripping, and NGC~628
vertical gravity give favorable results against the tested QUMOND
prescriptions or NFW populations. High-redshift kinematics remains a
tension. Established benchmarks assess recovery of acceleration, lensing,
acoustic spectra, expansion distances, cluster gravity, and tensor
propagation. Their evidential status is stated separately for each test.

\begingroup\small
\setlength{\tabcolsep}{4pt}
\renewcommand{\arraystretch}{1.10}
\begin{longtable}{@{}p{.06\textwidth}p{.24\textwidth}p{.56\textwidth}p{.06\textwidth}@{}}
\caption{All fifteen test families in the evidence register. Placement
identifies the role of a test; the status column states its present reach.
The register is exploratory, with the existing outcomes known.}\label{tab:fifteen}\\
\toprule ID & Test & Present result and limit & See\\\midrule
\endfirsthead
\toprule ID & Test & Present result and limit & See\\\midrule
\endhead
\multicolumn{4}{@{}l}{\textbf{I. Distinctive predictions and discriminating tests}}\\[3pt]
S01 & High-redshift kinematics & Reference line projection agrees; galaxy catalogue residuals retain a high-redshift tension. & \ref{sec:redshift}\\
S02 & CMB--galaxy normalization & Exponential-fit $a_0=1.25883\times10^{-10}\,\mathrm{m\,s^{-2}}$; standard linear-depletion value $1.25005\times10^{-10}\,\mathrm{m\,s^{-2}}$; conditional mass and survival factors. & \ref{sec:crossnorm}\\
S10 & Binaries and Solar System & Cassini null benchmark is compatible; a Gaia catalogue likelihood remains open. & \ref{sec:binaries}\\
S11 & $G$ and lepton masses & Close selected-formula values; finite routing tests leave discovery bias unquantified. & \ref{sec:constants}\\
S13 & Milky Way, LMC, streams, tilt & LMC mass agreement and repeated tilt realizations; Galactic mass and a joint stream fit remain unresolved. & \ref{sec:lmc}, \ref{sec:halo}\\
S14 & Satellites and tidal dwarfs & Crater II favors stock within the tested grids; Antlia 2 and odds depend strongly on model error. & \ref{sec:crater}\\
S15 & Old-disk vertical gravity & NGC~628 gives $\chi^2=2.04$ for five dispersions; a broader tracer and force comparison is needed. & \ref{sec:olddisk}\\
\midrule
\multicolumn{4}{@{}l}{\textbf{II. Recovery of established phenomenology}}\\[3pt]
S03 & Local RAR and BTFR & Broad SPARC diagnostic and small BTFR scatter; the score is sensitive to velocity errors. & \ref{sec:localrar}\\
S04 & Galaxy weak lensing & Updated fifteen-bin diagonal score 10.06; reduction sensitivity and joint covariance remain. & \ref{sec:galaxylensing}\\
S05 & CMB acoustic peaks & Planck-based matched fit and ACT comparison recover similar quality to the reference cosmology. & \ref{sec:cmb}\\
S06 & Growth and Weyl power & ACT lensing score 16.21 over ten bins; DESI full-shape and nonlinear growth tests remain open. & \ref{sec:cmb}\\
S07 & Dark-energy history and distances & Matched BAO/SN fit with exponential allocation; the feedback-resolved retention occupation gives $\Omega_\Lambda=0.709\pm0.02$ against $0.707$ with no input adjusted to it, and the frozen history fits with four parameters at $\Delta\chi^2=+1.07$. & \ref{sec:distances}\\
S08 & Relaxed and colliding clusters & Positive stock targets and conserved encounters; outlier mechanisms and lens-map fits remain conditional. & \ref{sec:clusters}\\
S09 & Slip and $E_G$ & Transfer predictions are compatible with two published estimates; a joint window-convolved fit is needed. & \ref{sec:eg}\\
S12 & Waves and Lorentz tests & Luminal tensor branch is compatible with the polarization benchmark; operator-level Lorentz tests remain open. & \ref{sec:tensor}\\
\bottomrule
\end{longtable}
\endgroup
\source{U1; CL1--CL3; D1--D2; T6}

The following calculations retain their individual likelihoods,
selection effects, and error models. A simulation realization, a fitted
relation, a conditional normalization, and an observation predicted
before measurement have different evidential roles. Lean checks the
stated mathematical implications; observed agreement is assessed
separately for each test.

\paragraph{Relation to the theory and CDT accounts.}
The effective astrophysical calculations use \emph{Entropic Scalar EFT}.
The companion CDT account distinguishes the electron-anchored scale
$G_*$, the selected bare coefficient $G_0\simeq1.03002G_*$, and an
interacting response $G_R$. The measured Newton constant remains the
effective input here. Matter response and the continuum phase of the
$\alpha=7$ CDT branch retain their own tests. \source{T1--T5}

\section{The common microscopic coefficients}
\label{sec:micro}

\subsection{Capacity, matter, and an observable gravitational field}

The starting hypothesis assigns finite entanglement capacity to a
tetrahedral substrate. Matter consists of persistent defects that commit
some of this capacity. Mass measures the commitment in energy units.
The ordinary gravitational branch identifies a capacity deficit with the
weak-field metric potential. In the manuscript's convention,
\begin{equation}
\frac{\delta S}{S_\infty}=-\frac{2\Phi}{c^2}.
\end{equation}
For a negative Newtonian potential, the corresponding capacity deficit is
positive. The ordinary metric branch uses the measured Newton constant. Varying
the specified two-potential action independently in $\Phi$ and $\Psi$
gives
\begin{equation}
\nabla^2\Psi=4\pi G\rho,\qquad
\nabla^2(\Phi-\Psi)=0.
\end{equation}
For smooth isolated sources with $\Phi-\Psi\to0$ at infinity, the
maximum principle gives $\Phi=\Psi$. These steps have kernel-checked
proofs, connecting the reduced action to the Newtonian force and the
leading no-slip metric. The physical selection of that action and its
normalization retain their stated premises. \source{T4}

A stellar velocity, a stream deflection, and a gravitational lens measure
the field generated by the effective stress-energy. They do not directly
identify the microscopic material carrying that stress. The galaxy runs
represent retained capacity by collisionless particles. Those numerical
particles sample a positive mass distribution; the theory interprets their
mass as substrate commitment. Thus the model contains an additional
gravitating component even though it proposes no new species of dark-matter
particle. Its physical origin and relation to baryons are the claims under
test. \source{T1, sections 1, 15, 20 and O.3}

\subsection{Why the selected face has seven transmitted and nine retained states}

Each side of a shared face carries a primitive spin-$3/2$ representation,
which has four states. Its paired space has dimension sixteen. Angular
momentum addition gives
\begin{equation}
V_{3/2}\otimes V_{3/2}=V_0\oplus V_1\oplus V_2\oplus V_3,
\qquad 16=1+3+5+7.
\end{equation}
The selected positive matching operator transmits the maximal-spin sector
$V_3$. Its orthogonal complement has dimensions $1+3+5=9$ and supplies the
retained marked register. The present/history response has the same
representation content, $V_1\otimes V_1=V_0\oplus V_1\oplus V_2$.
An explicit equivariant map identifies the marked register with the
retained complement. Its uniqueness is proved up to one scalar factor
on each spin block; equivariant unitary freedom consists of block phases.
The result fixes the allowed identification within the selected vertex.
\source{T4}

The numbers describe dimensions of state spaces. They do not say that
seven sixteenths of the physical energy passes and nine sixteenths becomes
dark energy. An energy fraction requires an energy operator and an
identified state. The release fraction derived below is about three
percent, not $9/16$.

There are four distinguishable face ports. Single-copy occupancy makes
their seven-letter assignments injective. Including the two cell
orientations gives
\begin{equation}
N_{\rm states}=2P(7,4)=2\,(7\cdot6\cdot5\cdot4)=1680.
\end{equation}
Uniform weights would give entropy $\ln1680=7.4265490724$. The theory
weights these states by their closure defect, producing a slightly smaller
entropy. \source{T1, sections 5--6 and B.2}

\subsection{The finite sum that fixes the entropy}

Let $m_i\in\{-3,-2,-1,0,1,2,3\}$ be the four distinct face labels.
The tetrahedral face normals satisfy
$\widehat n_i\cdot\widehat n_j=-1/3$ for $i\ne j$. On the product-state
boundary realization the expectation of the squared closure operator is
\begin{equation}
K^2=48-\frac13\left[\left(\sum_i m_i\right)^2-\sum_i m_i^2\right].
\label{eq:K}
\end{equation}
The constant $48=4j(j+1)$ uses the transmitted spin $j=3$.
The Gibbs weights and their partition sum are
\begin{equation}
p_\eta(b)=\frac{e^{-\eta K^2(b)}}{Z(\eta)},\qquad
Z(\eta)=\sum_{b=1}^{1680}e^{-\eta K^2(b)}.
\end{equation}
Three closure components supply the determinant factor $\eta^{3/2}$.
Stationarity of the normalized closure evidence gives
\begin{equation}
0=\frac{d}{d\eta}\left[\ln Z+\frac32\ln\eta\right]
=-\langle K^2\rangle_\eta+\frac{3}{2\eta}.
\end{equation}
Direct enumeration and a one-dimensional root solve yield
\begin{align}
\eta_*&=0.0298668443935, & \langle K^2\rangle_*&=50.2229154254,\\
g&=-\sum_b p_*(b)\ln p_*(b)=7.41980002357.
\label{eq:g}
\end{align}
Finite enumeration and the closure root have kernel-checked proofs and
certified numerical bounds. Independent floating-point evaluation supplies
the additional displayed digits; the verification records distinguish the
two forms of calculation. \source{T4}
The word ``fixed'' means fixed after the ensemble, matching rule,
closure invariant, and evidence measure have been specified. It does not
claim that finite capacity uniquely determines that microscopic model.

\subsection{Two coefficients inherited by the astrophysical branches}

The transverse galactic branch takes its response cell to be the tilt
mode of the transverse frame, with inertia lever $\ell=g\,c/H$. The
horizon temperature $\hbar H/2\pi$ then gives the dimensionless coefficient
\begin{equation}
\epsilon=\frac{g}{4\pi^2}=0.1879457302,
\qquad \epsilon^{-1}=5.320684854.
\label{eq:eps}
\end{equation}
and the acceleration scale
\begin{equation}
a_0(z)=\epsilon cH(z).
\label{eq:a0}
\end{equation}
At $H_0=67.4\,\mathrm{km\,s^{-1}\,Mpc^{-1}}$,
$a_0(0)=1.23072924\times10^{-10}\,\mathrm{m\,s^{-2}}$.
The standard linear-depletion fit has $H_0=68.458\,\mathrm{km\,s^{-1}\,Mpc^{-1}}$
and $a_0(0)=1.25005\times10^{-10}\,\mathrm{m\,s^{-2}}$.
The exponential-allocation fit with massless neutrinos, linear lensing and
a low-redshift supernova offset has $H_0=68.93885586416\,\mathrm{km\,s^{-1}\,Mpc^{-1}}$
and $a_0(0)=1.25883\times10^{-10}\,\mathrm{m\,s^{-2}}$.
The galaxy calculations specify the reference or fitted normalization used
in each comparison.
This coefficient is obtained by evaluating a fixed prescription.
The tilt-mode identification, its horizon thermalization, and the lever
are the physical matching assumptions of that prescription; the lever
coefficient follows from the register enthalpy and a one-nat oriented record
read at the Unruh temperature of the projected field, given that reading. \source{T1, section 15}

In the cosmological branch, the committed state has mean energy
$\varepsilon_K\langle K^2\rangle_*$, while the three normalized record
coordinates carry mean energy $3\varepsilon_\xi/2$. Their ratio is
\begin{equation}
\eta=\frac{\varepsilon_\xi}{\varepsilon_K}
\frac{3}{2\langle K^2\rangle_*}.
\label{eq:etaenergy}
\end{equation}
The common-conversion choice $\varepsilon_\xi=\varepsilon_K$ makes
$\eta=\eta_*$. This identification also uses the manuscript's post-refresh
energy accounting. The exact closure identity and the physical equality
of conversion scales have different logical roles. \source{T1, O.4 and O.13}

Here $g$ is the sharing entropy denoted $g_{\rm share}$ in the CDT
account and $g_{\rm share,eff}$ in Astra. The galactic $\epsilon$ is
distinct from the electron rest gap $\epsilon_e\simeq4.18546\times10^{-23}$.

The same pair $(g,\eta_*)$ therefore enters the galactic scale, the initial
commitment budget, and the release fraction. The observational calculations
test those connections together with the effective dynamics that takes
them from a cell to a galaxy or an expanding universe.
